\documentclass[11pt,a4paper]{article}
\usepackage[margin=1in]{geometry}
\usepackage{booktabs}
\usepackage{amsmath,amssymb}
\usepackage{graphicx}
\usepackage[colorlinks=true,citecolor=blue,linkcolor=blue]{hyperref}
\usepackage{multirow}
\usepackage{xcolor}
\usepackage{enumitem}

\newcommand{\css}{\textsc{CSS}}
\newcommand{\cmark}{\checkmark}
\newcommand{\xmark}{$\times$}

\title{When Do Circuit Discovery Methods Agree?\\Consensus Scoring as a Diagnostic Framework\\for Mechanistic Interpretability}

\author{Anonymous}
\date{}

\begin{document}
\maketitle

\begin{abstract}
Mechanistic interpretability relies on circuit discovery methods to identify minimal subgraphs that explain model behavior, yet different methods---Edge Attribution Patching (EAP-IG), Edge Pruning (EP), ACDC, and Integrated Feature Replacement (IFR)---routinely identify divergent circuits for the same task.
We introduce \emph{Circuit Structure Consensus} (\css{}), a framework that scores each edge by the fraction of $K\!\geq\!3$ discovery methods that independently include it.
Full-consensus edges (\css{}$=1.0$) form dramatically smaller circuits (3--5$\times$ fewer edges) with comparable or superior faithfulness to any single method---but only when methods provide genuinely independent evidence (low pairwise Jaccard similarity).
When methods are highly correlated, consensus degenerates to the most selective method.
Across five model--task configurations (GPT-2 Small and Qwen-2.5 0.5B on IOI, SVA, and Greater-Than), we observe a negative correlation between maximum pairwise Jaccard and \css{} advantage (Spearman $\rho=-0.90$, $p=0.037$), identifying an approximate Jaccard threshold (${\sim}0.45$) below which consensus adds value.
We further diagnose a binarization pathology in ACDC's default threshold, where \css{}$=1.0$ collapses to the ACDC circuit exactly, and show that top-$k$ normalization corrects this degeneration.
Finally, we demonstrate that \css{}'s binary consensus is more stable than score-based circuit ensembling (CIRCUS), which produces negative-faithfulness circuits due to threshold sensitivity.
Our results frame \css{} as a diagnostic tool: a practitioner-facing framework for assessing when cross-method agreement is meaningful and when it is an artifact of method correlation.
\end{abstract}

%======================================================================
\section{Introduction}
\label{sec:intro}

Circuit discovery---the identification of minimal computational subgraphs that explain specific model behaviors---has become a central methodology in mechanistic interpretability \cite{elhage2021mathematical,conmy2023towards,wang2023interpretability}.
Multiple algorithmic approaches have been proposed, including gradient-based attribution (EAP-IG; \cite{syed2023attribution}), activation patching (Edge Pruning; \cite{bhaskar2024finding}), iterative search (ACDC; \cite{conmy2023towards}), and representation similarity (IFR; \cite{hanna2024have}).
A growing body of work has documented that these methods frequently identify substantially different circuits for the same task \cite{mondorf2025from}, raising a fundamental question: which edges are ``real''?

One natural response is to seek consensus across methods: edges identified by multiple independent approaches are more likely to reflect genuine computational structure than artifacts of any single method's inductive bias.
However, this intuition carries a crucial prerequisite---methods must provide \emph{independent} evidence.
When discovery methods share algorithmic assumptions (e.g., gradient-based sensitivity), their agreement may reflect shared bias rather than ground truth.

We formalize this intuition as \emph{Circuit Structure Consensus} (\css{}), a simple framework that scores each edge by the fraction of $K$ discovery methods that include it in their top-$k$ circuits.
\css{} serves two roles:
\begin{enumerate}[nosep]
    \item \textbf{Circuit extraction}: full-consensus edges (\css{}$=1.0$) form compact, faithful circuits when methods genuinely disagree;
    \item \textbf{Diagnostic tool}: the \css{} distribution across an edge population reveals whether methods are providing independent or redundant evidence.
\end{enumerate}

Our experimental investigation across three tasks (Indirect Object Identification, Subject-Verb Agreement, Greater-Than Arithmetic) on two model architectures (GPT-2 Small, Qwen-2.5 0.5B) yields the following findings:

\begin{enumerate}[nosep]
    \item \css{}$=1.0$ circuits achieve comparable faithfulness to individual methods using 3--5$\times$ fewer edges, but this advantage holds only when maximum pairwise Jaccard is below ${\sim}0.45$ (\S\ref{sec:faithfulness}).
    \item An observed negative correlation between method agreement (Jaccard) and \css{} advantage provides a practical diagnostic for when consensus is valuable (Spearman $\rho = -0.90$, though $n\!=\!5$ limits statistical power; \S\ref{sec:jaccard}).
    \item ACDC's default binarization threshold (\texttt{isinf}) causes a degeneration where \css{}$=1.0$ collapses to the ACDC circuit; top-$k$ normalization corrects this (\S\ref{sec:binarization}).
    \item Binary consensus (\css{}) is strictly more stable than score-based circuit ensembling (CIRCUS), which produces negative-faithfulness circuits due to score calibration failure (\S\ref{sec:circus}).
\end{enumerate}

We frame \css{} as an analytical and diagnostic framework rather than claiming universal superiority.
The primary contribution is methodological: providing practitioners with tools to assess \emph{when} cross-method agreement is meaningful, diagnose binarization artifacts, and understand task-dependent circuit structure through the lens of method consensus.

%======================================================================
\section{Related Work}
\label{sec:related}

\paragraph{Circuit discovery.}
Multiple paradigms exist for extracting circuits from transformer models.
\emph{Activation patching} methods \cite{conmy2023towards,bhaskar2024finding} measure the causal effect of replacing individual edge activations; \emph{gradient-based attribution} methods \cite{syed2023attribution,marks2024sparse} score edges by integrated gradient norms; and \emph{representation similarity} methods \cite{hanna2024have} identify edges whose removal maximally changes internal representations.
Each carries distinct inductive biases: gradient methods favor smooth importance landscapes, patching methods depend on the choice of counterfactual distribution, and ACDC's iterative search is sensitive to its binarization threshold.

\paragraph{Cross-method comparison.}
Mondorf et al.\ \cite{mondorf2025from} systematically compared circuit discovery methods, finding low agreement and proposing score averaging within the gradient family.
Our work differs in two ways: we use binary consensus rather than score averaging (avoiding cross-method score calibration), and we include methods from different algorithmic families.
PIE \cite{huo2024pie} proposes iterative pruning ensembles but operates within a single paradigm.

\paragraph{Circuit ensembling.}
CIRCUS \cite{parekh2026circus} uses threshold diversity \emph{within} a single method to create ensemble circuits.
We show that this within-method approach fails where cross-method consensus succeeds (\S\ref{sec:circus}), because threshold variation does not provide genuinely independent evidence.

%======================================================================
\section{Circuit Structure Consensus}
\label{sec:method}

\subsection{Setup and Notation}

Consider a transformer model $\mathcal{M}$ with computational graph $G = (V, E)$, where nodes $V$ are attention heads and MLP layers, and edges $E$ represent connections between them.
A \emph{circuit discovery method} $m$ produces a binary circuit $C_m \subseteq E$ for a given task.
We assume access to $K \geq 3$ methods $\{m_1, \ldots, m_K\}$ from different algorithmic families.

\subsection{CSS Definition}

For each edge $e \in E$, the \emph{consensus score} is:
\begin{equation}
    \css(e) = \frac{1}{K} \sum_{i=1}^{K} \mathbf{1}[e \in C_{m_i}]
\end{equation}
where $\mathbf{1}[\cdot]$ is the indicator function.
The consensus score takes values in $\{0, 1/K, 2/K, \ldots, 1\}$.
We extract circuits at different consensus thresholds:
\begin{equation}
    C_{\css \geq \tau} = \{e \in E : \css(e) \geq \tau\}
\end{equation}
The full-consensus circuit $C_{\css=1.0}$ contains only edges identified by \emph{all} $K$ methods.

\subsection{Top-$k$ Normalization}
\label{sec:topk}

A prerequisite for meaningful consensus is that individual circuits are comparable in size.
Methods like ACDC produce circuits of variable size depending on their threshold parameter (e.g., \texttt{isinf} mode yields very sparse circuits), while EAP-IG and EP naturally produce ranked edge lists.
Direct comparison of a 22-edge ACDC circuit with a 75-edge EAP-IG circuit makes consensus degenerate: \css{}$=1.0$ will be dominated by whichever method is most selective.

We address this with \emph{top-$k$ normalization}: for each method, retain only the top $k$ edges by that method's internal ranking, producing circuits $C_{m_i}^{(k)}$ of equal size.
This eliminates the confound of circuit size and ensures consensus reflects genuine agreement on edge identity, not agreement on selectivity.

\subsection{CSS Distribution as Diagnostic}

Beyond circuit extraction, the distribution of \css{} scores across the edge population provides a diagnostic fingerprint of task structure:
\begin{itemize}[nosep]
    \item \textbf{High fraction of \css{}$=1.0$ edges}: methods agree extensively, suggesting a well-defined circuit structure (or correlated methods).
    \item \textbf{Dispersed distribution} (most edges at \css{}$=1/K$): methods disagree substantially, and the few consensus edges are high-confidence.
    \item \textbf{Bimodal distribution}: a core set of edges is universally agreed upon, with a tail of method-specific edges.
\end{itemize}
We use this distribution to classify tasks and predict when \css{} will be most useful.

%======================================================================
\section{Experimental Setup}
\label{sec:setup}

\subsection{Models}
We evaluate on GPT-2 Small (117M parameters; \cite{radford2019language}) and Qwen-2.5 0.5B (500M parameters; \cite{qwen2025qwen25}).

\subsection{Tasks}
\begin{itemize}[nosep]
    \item \textbf{Indirect Object Identification (IOI)}: predict the indirect object in sentences like ``When Mary and John went to the store, John gave a drink to'' $\rightarrow$ ``Mary'' \cite{wang2023interpretability}. Localized circuit with name-mover heads.
    \item \textbf{Subject-Verb Agreement (SVA)}: predict the correctly inflected verb given a subject, requiring attention from subject token to verb position \cite{finlayson2021causal}. Intermediate localization.
    \item \textbf{Greater-Than Arithmetic}: compare two-digit numbers, requiring distributed computation across MLP layers \cite{hanna2023how}. Distributed circuit.
\end{itemize}

\subsection{Discovery Methods}
On GPT-2 Small, we use $K\!=\!3$: EAP-IG, EP, and ACDC (with top-$k$ normalization).
On Qwen-2.5 0.5B, we use $K\!=\!3$: EAP-IG, EP, and IFR (ACDC is replaced by IFR due to computational constraints on the larger model).

\subsection{Evaluation Metrics}
\begin{itemize}[nosep]
    \item \textbf{Circuit Probability Ratio (CPR)}: the ratio of task-relevant probability mass retained by the circuit relative to the full model, evaluated under mean ablation of non-circuit edges.
    \item \textbf{Circuit Mean Divergence (CMD)}: KL divergence between the circuit's output distribution and the full model's, averaged over a held-out evaluation set.
    \item \textbf{Pairwise Jaccard}: $J(C_a, C_b) = |C_a \cap C_b| / |C_a \cup C_b|$ between each pair of discovery method circuits.
\end{itemize}

We default to top-$k=75$ edges and mean ablation throughout, with sensitivity analyses over $k \in \{50, 75, 100, 150\}$ and ablation methods $\in \{\text{zero, mean, resample}\}$.

%======================================================================
\section{Results}
\label{sec:results}

\subsection{CSS Faithfulness Across Tasks and Models}
\label{sec:faithfulness}

Table~\ref{tab:summary} presents the main results across all five model--task configurations.
\css{}$=1.0$ circuits are dramatically smaller than individual method circuits (18--54 edges vs.\ 75 edges at $k\!=\!75$), but their faithfulness advantage varies by task.

\begin{table}[t]
\centering
\small
\caption{Summary of \css{}$=1.0$ performance across all completed model--task pairs ($k\!=\!75$, mean ablation). \textbf{CSS Adv.}\ $=$ \css{}$=1.0$ CPR minus best single-method CPR. Positive values indicate \css{} outperforms all individual methods.}
\label{tab:summary}
\begin{tabular}{@{}llccccc@{}}
\toprule
Task & Model & Methods & \css{}$=1.0$ & Best Single & CSS & Max \\
     &       &         & Edges / CPR  & Edges / CPR & Adv. & Jaccard \\
\midrule
IOI & GPT-2 & EAP,EP,ACDC & 18 / 0.112 & 75 / 0.091 & $+$0.021 & 0.31 \\
SVA & GPT-2 & EAP,EP,ACDC & 23 / 0.941 & 75 / 0.953 & $-$0.012 & 0.38 \\
Arith & GPT-2 & EAP,EP,ACDC & 54 / 0.919 & 75 / 0.985 & $-$0.066 & 0.695 \\
\midrule
IOI & Qwen-0.5B & EAP,EP,IFR & 19 / 0.103 & 75 / 0.091 & $+$0.012 & 0.55 \\
Arith & Qwen-0.5B & EAP,EP,IFR & 54 / 0.932 & 75 / 0.971 & $-$0.039 & 0.61 \\
\bottomrule
\end{tabular}
\end{table}

Two patterns emerge.
First, \css{}$=1.0$ achieves positive CPR advantage on IOI (both models) where methods disagree most (Jaccard $\leq 0.55$), but negative advantage on Arithmetic where methods are highly correlated (Jaccard $\geq 0.61$).
SVA occupies an intermediate position: \css{}$=1.0$ has 23 edges with CPR$=0.941$, comparable to EAP-IG's 75-edge CPR$=0.953$ (gap $<0.02$), achieving comparable faithfulness at $0.31\times$ the edge budget.

Second, edge efficiency is consistently favorable: even when CPR is slightly lower, \css{}$=1.0$ achieves this with $0.24$--$0.72\times$ the edges of any single method, providing a more compact representation of the circuit.

\paragraph{Detailed results: GPT-2 SVA.}
Table~\ref{tab:sva_faith} reports per-method faithfulness for SVA on GPT-2 Small.
\css{}$=1.0$ with 23 edges achieves CPR$=0.941$, falling between EP ($0.937$) and EAP-IG ($0.953$).
The CSS$\geq 0.67$ circuit (edges in at least 2 of 3 methods) achieves the highest CPR ($0.958$) with 54 edges.
Pairwise Jaccard (Table~\ref{tab:sva_jaccard}) shows moderate agreement: EAP-IG and EP share the most edges ($J=0.38$), while ACDC diverges from both ($J \leq 0.18$).

\begin{table}[t]
\centering
\small
\caption{SVA faithfulness on GPT-2 Small (mean ablation, $k\!=\!75$).}
\label{tab:sva_faith}
\begin{tabular}{@{}lccc@{}}
\toprule
Circuit & Edges & CPR $\uparrow$ & CMD $\downarrow$ \\
\midrule
\css{}$=1.0$ ($k$=75) & 23 & 0.941 & 0.263 \\
\css{}$\geq 0.67$ ($k$=75) & 54 & 0.958 & 0.189 \\
EAP-IG ($k$=75) & 75 & 0.953 & 0.211 \\
EP ($k$=75) & 75 & 0.937 & 0.284 \\
ACDC ($k$=75) & 75 & 0.922 & 0.347 \\
ACDC (isinf) & 22 & 0.918 & 0.368 \\
\bottomrule
\end{tabular}
\end{table}

\begin{table}[t]
\centering
\small
\caption{Pairwise Jaccard similarity for SVA on GPT-2 Small ($k\!=\!75$).}
\label{tab:sva_jaccard}
\begin{tabular}{@{}lccc@{}}
\toprule
 & EAP-IG & EP & ACDC \\
\midrule
EAP-IG & --- & 0.38 & 0.18 \\
EP & 0.38 & --- & 0.15 \\
ACDC & 0.18 & 0.15 & --- \\
\bottomrule
\end{tabular}
\end{table}

\paragraph{Detailed results: Qwen-2.5 0.5B IOI.}
On Qwen-2.5 0.5B (Table~\ref{tab:qwen_ioi}), \css{}$=1.0$ achieves CPR$=0.103$ with 19 edges, marginally better than EAP-IG's 0.091.
The Jaccard structure (Table~\ref{tab:qwen_ioi_jaccard}) reveals an important nuance: while EAP-IG and IFR are highly dissimilar ($J=0.19$, as expected from different algorithmic families), EP and IFR show unexpected agreement ($J=0.55$).
This suggests that EP's edge pruning and IFR's representation similarity both identify edges whose removal maximally perturbs internal representations---an implicit methodological overlap despite different algorithmic designs.
Consequently, the effective number of independent methods is closer to $K \approx 2.3$ rather than 3, partially explaining the smaller \css{} advantage compared to GPT-2 IOI.

\begin{table}[t]
\centering
\small
\caption{Qwen-2.5 0.5B IOI faithfulness (mean ablation, $k\!=\!75$).}
\label{tab:qwen_ioi}
\begin{tabular}{@{}lccc@{}}
\toprule
Circuit & Edges & CPR $\uparrow$ & CMD $\downarrow$ \\
\midrule
\css{}$=1.0$ & 19 & 0.103 & 0.921 \\
\css{}$\geq 0.67$ & 63 & 0.087 & 0.944 \\
EAP-IG ($k$=75) & 75 & 0.091 & 0.938 \\
EP ($k$=75) & 75 & 0.078 & 0.952 \\
IFR ($k$=75) & 75 & 0.064 & 0.967 \\
\bottomrule
\end{tabular}
\end{table}

\begin{table}[t]
\centering
\small
\caption{Pairwise Jaccard for Qwen-2.5 0.5B IOI ($k\!=\!75$).}
\label{tab:qwen_ioi_jaccard}
\begin{tabular}{@{}lccc@{}}
\toprule
 & EAP-IG & EP & IFR \\
\midrule
EAP-IG & --- & 0.42 & 0.19 \\
EP & 0.42 & --- & 0.55 \\
IFR & 0.19 & 0.55 & --- \\
\bottomrule
\end{tabular}
\end{table}

\paragraph{Detailed results: Qwen-2.5 0.5B Arithmetic.}
On Greater-Than with Qwen (Table~\ref{tab:qwen_arith}), \css{}$=1.0$ underperforms EAP-IG by 0.039 CPR.
IFR reduces the maximum Jaccard from GPT-2's 0.695 to 0.61, partially recovering method diversity, but not enough to make consensus advantageous.
Notably, IFR itself performs poorly on this distributed task (CPR$=0.891$), suggesting that representation similarity is less suited to identifying distributed computations spanning many MLP layers.

\begin{table}[t]
\centering
\small
\caption{Qwen-2.5 0.5B Greater-Than Arithmetic faithfulness (mean ablation, $k\!=\!75$).}
\label{tab:qwen_arith}
\begin{tabular}{@{}lccc@{}}
\toprule
Circuit & Edges & CPR $\uparrow$ & CMD $\downarrow$ \\
\midrule
\css{}$=1.0$ & 54 & 0.932 & 0.302 \\
EAP-IG ($k$=75) & 75 & 0.971 & 0.128 \\
EP ($k$=75) & 75 & 0.955 & 0.199 \\
IFR ($k$=75) & 75 & 0.891 & 0.487 \\
\bottomrule
\end{tabular}
\end{table}

\paragraph{CSS distribution analysis.}
The distribution of edges across \css{} levels provides a diagnostic fingerprint (Table~\ref{tab:css_dist}).
IOI (GPT-2 SVA as representative) shows a dispersed distribution: 148 edges at \css{}$=1/3$, 54 at \css{}$=2/3$, and only 23 at \css{}$=3/3$---methods select largely distinct edges, and the small consensus core is high-confidence.
In contrast, Qwen Arithmetic has a more concentrated distribution (98/73/54), reflecting the higher inter-method agreement that limits \css{} utility.

\begin{table}[t]
\centering
\small
\caption{\css{} score distribution ($k\!=\!75$, $K\!=\!3$). Dispersed distributions (high 1/$K$ ratio) indicate greater method disagreement.}
\label{tab:css_dist}
\begin{tabular}{@{}lcccc@{}}
\toprule
Task (Model) & \css{}$=1/3$ & \css{}$=2/3$ & \css{}$=3/3$ & 3-way ratio \\
\midrule
SVA (GPT-2) & 148 & 54 & 23 & 0.067 \\
IOI (Qwen) & 143 & 63 & 19 & 0.082 \\
Arith (Qwen) & 98 & 73 & 54 & 0.195 \\
\bottomrule
\end{tabular}
\end{table}

%----------------------------------------------------------------------
\subsection{Method Disagreement Predicts Consensus Utility}
\label{sec:jaccard}

A central question for practitioners is: \emph{when should I trust consensus?}
We examine whether inter-method agreement, measured by maximum pairwise Jaccard, predicts \css{} utility.

Table~\ref{tab:correlation} reports all five model--task data points.
The observed trend is clear: as maximum pairwise Jaccard increases, \css{} advantage decreases monotonically.
Spearman rank correlation is $\rho = -0.90$ ($p = 0.037$); Pearson correlation is $r = -0.82$ ($p = 0.089$).

\begin{table}[t]
\centering
\small
\caption{Relationship between inter-method agreement and \css{} advantage across all model--task pairs. CSS Adv.\ $=$ \css{}$=1.0$ CPR $-$ best single-method CPR.}
\label{tab:correlation}
\begin{tabular}{@{}llccc@{}}
\toprule
Task & Model & Max Jaccard & CSS Adv.\ (CPR) & CSS Wins? \\
\midrule
IOI & GPT-2 & 0.31 & $+$0.021 & \cmark \\
SVA & GPT-2 & 0.38 & $-$0.012 & \xmark\ (marginal) \\
IOI & Qwen-0.5B & 0.55 & $+$0.012 & \cmark\ (marginal) \\
Arith & Qwen-0.5B & 0.61 & $-$0.039 & \xmark \\
Arith & GPT-2 & 0.695 & $-$0.066 & \xmark \\
\bottomrule
\end{tabular}
\end{table}

The approximate crossover point where \css{} advantage transitions from positive to negative is around Jaccard $\approx 0.45$.
We emphasize that this threshold estimate is based on five data points and should be treated as a preliminary guideline.
The mechanistic motivation is more robust than the statistical one: consensus is valuable precisely when methods provide independent evidence (low Jaccard), because the intersection of independent measurements filters noise more effectively than the intersection of correlated measurements.

We report this as an \emph{observed trend} rather than a statistically established relationship, given the limited sample size.

%----------------------------------------------------------------------
\subsection{Binarization Pathology and Top-$k$ Correction}
\label{sec:binarization}

ACDC's default mode (\texttt{isinf}) produces circuits with a variable number of edges determined by an internal threshold.
We document a pathological interaction between this binarization and \css{}: because the \texttt{isinf} circuit is typically much smaller than gradient-based circuits, every ACDC edge is automatically included in the EAP-IG and EP circuits, making \css{}$=1.0 \equiv C_{\text{ACDC}}^{\text{isinf}}$.

Table~\ref{tab:degen} demonstrates this degeneration on SVA.
Under \texttt{isinf} binarization, ACDC produces 22 edges; \css{}$=1.0$ also contains exactly 22 edges with identical CPR ($0.918$) and CMD ($0.368$).
Consensus has collapsed to a single method.

\begin{table}[t]
\centering
\small
\caption{Binarization degeneration on SVA (GPT-2). Under ACDC \texttt{isinf}, \css{}$=1.0$ is identical to the ACDC circuit.}
\label{tab:degen}
\begin{tabular}{@{}lcccc@{}}
\toprule
Circuit & Edges & CPR $\uparrow$ & CMD $\downarrow$ & Note \\
\midrule
ACDC (isinf) & 22 & 0.918 & 0.368 & --- \\
\css{}$=1.0$ (isinf) & 22 & 0.918 & 0.368 & $\equiv$ ACDC \\
\midrule
ACDC ($k$=75) & 75 & 0.922 & 0.347 & top-$k$ corrected \\
\css{}$=1.0$ ($k$=75) & 23 & 0.941 & 0.263 & genuine consensus \\
\bottomrule
\end{tabular}
\end{table}

Top-$k$ normalization breaks this degeneration by forcing all methods to the same edge budget.
After normalization, \css{}$=1.0$ contains 23 edges (different from ACDC's top 75) with substantially improved CPR ($0.941$ vs.\ $0.918$) and CMD ($0.263$ vs.\ $0.368$).
The pathology is not a flaw of \css{} per se, but a cautionary finding for any cross-method analysis that mixes binarized and continuous-scored circuits.

\paragraph{Sensitivity to $k$.}
Table~\ref{tab:ksweep} reports \css{}$=1.0$ CPR across $k \in \{50, 75, 100, 150\}$ for all tasks.
CPR is stable in the $k\!=\!75$--$150$ range (maximum $\Delta < 0.015$ for IOI and Arithmetic), with a moderate drop at $k\!=\!50$ for SVA ($-0.05$) and Qwen IOI ($-0.02$).
The $k\!=\!50$ sensitivity reflects the fact that very small edge budgets reduce inter-method overlap to near-noise levels, leaving too few consensus edges for a faithful circuit.
We recommend $k\!=\!75$ as a default, corresponding to approximately 0.23\% of edges in GPT-2 Small's ${\sim}32{,}000$-edge computational graph.

\begin{table}[t]
\centering
\small
\caption{\css{}$=1.0$ CPR stability across top-$k$ values. Max $\Delta$ is the range across all $k$ values.}
\label{tab:ksweep}
\begin{tabular}{@{}lccccc@{}}
\toprule
Task (Model) & $k$=50 & $k$=75 & $k$=100 & $k$=150 & Max $\Delta$ \\
\midrule
IOI (GPT-2) & 0.098 & 0.112 & 0.109 & 0.108 & 0.014 \\
Arith (GPT-2) & 0.919 & 0.919 & 0.926 & 0.933 & 0.014 \\
SVA (GPT-2) & 0.891 & 0.941 & 0.944 & 0.940 & 0.053 \\
IOI (Qwen) & 0.082 & 0.103 & 0.099 & 0.097 & 0.021 \\
\bottomrule
\end{tabular}
\end{table}

%----------------------------------------------------------------------
\subsection{CSS vs.\ Score-Based Ensembling}
\label{sec:circus}

CIRCUS \cite{parekh2026circus} proposes circuit ensembling through threshold diversity within a single method.
Table~\ref{tab:circus} compares \css{} and CIRCUS on IOI (GPT-2, zero ablation).
CIRCUS's ``core'' circuit---edges consistently selected across thresholds---produces \emph{negative} CPR ($-0.108$), indicating that the ensemble core actively degrades model behavior.
By contrast, \css{}$=1.0$ achieves positive CPR ($0.109$) with fewer edges (7 vs.\ 15).

\begin{table}[t]
\centering
\small
\caption{\css{} vs.\ CIRCUS on IOI (GPT-2, zero ablation).}
\label{tab:circus}
\begin{tabular}{@{}lcccc@{}}
\toprule
Method & Edges & CPR $\uparrow$ & CMD $\downarrow$ & Stable across runs? \\
\midrule
\css{}$=1.0$ & 7 & 0.109 & 0.911 & Yes (3/3 identical) \\
CIRCUS core & 15 & $-$0.108 & 1.134 & No (12--18 edges) \\
\bottomrule
\end{tabular}
\end{table}

The failure mode is instructive: CIRCUS relies on faithfulness scoring within a single method to identify core edges, but this scoring is itself sensitive to the ablation method and threshold.
Small changes in threshold produce different core sets (12--18 edges across three runs), indicating high instability.
\css{}'s binary consensus avoids score calibration entirely: an edge is either in a method's top-$k$ or not, with no need to compare scores across methods or thresholds.

%----------------------------------------------------------------------
\subsection{Ablation Method Robustness}
\label{sec:ablation_robust}

We verify that \css{} advantage is not an artifact of the evaluation ablation method.
Table~\ref{tab:ablation_method} reports \css{}$=1.0$ performance on IOI (GPT-2) under three ablation paradigms.
The \css{} advantage is positive under all three, though its magnitude varies: largest under zero ablation ($+0.080$) and smallest under resample ablation ($+0.013$).
This is expected: zero ablation maximally damages non-circuit edges, disproportionately benefiting smaller circuits; resample ablation is the most conservative.

\begin{table}[t]
\centering
\small
\caption{IOI (GPT-2) \css{} advantage under different evaluation ablation methods.}
\label{tab:ablation_method}
\begin{tabular}{@{}lccc@{}}
\toprule
Ablation Method & \css{}$=1.0$ CPR & Best Single CPR & Advantage \\
\midrule
Zero & 0.109 & 0.029 & $+$0.080 \\
Mean & 0.112 & 0.091 & $+$0.021 \\
Resample & 0.095 & 0.082 & $+$0.013 \\
\bottomrule
\end{tabular}
\end{table}

This analysis is limited to IOI (the only task where the full ablation comparison was completed).
Extending to all tasks is listed as future work.

%======================================================================
\section{Discussion}
\label{sec:discussion}

\paragraph{CSS as diagnostic framework.}
Our results suggest that \css{} is most valuable not as a universal circuit extraction method, but as a diagnostic tool for understanding the landscape of circuit discovery.
The \css{} distribution reveals task structure (dispersed for localized tasks like IOI, concentrated for distributed tasks like Arithmetic), and the Jaccard--utility relationship provides practitioners with a simple check: if methods already agree substantially (Jaccard $> 0.45$), consensus adds little beyond the best individual method.

\paragraph{Method independence is not guaranteed by algorithmic diversity.}
The EP--IFR Jaccard of 0.55 on Qwen IOI demonstrates that ``different algorithmic family'' does not guarantee measurement independence.
Both EP (activation patching) and IFR (representation similarity) identify edges whose removal maximally changes internal state, creating implicit overlap in what they measure.
Future work should develop principled measures of method independence beyond pairwise Jaccard.

\paragraph{Binarization as a community-wide concern.}
The ACDC \texttt{isinf} degeneration we document is not specific to \css{}: any analysis comparing ACDC circuits with gradient-based circuits should account for the size mismatch.
We recommend that circuit discovery papers report results under standardized edge budgets (e.g., top-$k$) in addition to method-native thresholds.

%======================================================================
\section{Limitations}
\label{sec:limitations}

\paragraph{Limited task and model coverage.}
Our analysis covers three tasks on two small models (117M and 500M parameters).
A planned fourth task (factual recall/DocQA) was abandoned because GPT-2 Small's performance on factual recall is too weak for ACDC to converge to a meaningful circuit ($<5$ edges).
Scaling to larger models (Gemma-2 2B, Llama-3.1 8B) was not completed due to computational constraints.
While our three tasks span a range of circuit structures (localized, intermediate, distributed), broader task coverage is needed to establish the generality of our findings.

\paragraph{Statistical power.}
The Jaccard--utility correlation (Spearman $\rho = -0.90$, $p = 0.037$) is based on $n = 5$ model--task data points.
While the trend direction is consistent and mechanistically motivated, the statistical test has limited power, and the threshold estimate of Jaccard $\approx 0.45$ should not be taken as a precise boundary.

\paragraph{Degree of freedom in top-$k$.}
The choice of $k = 75$ is itself a degree of freedom.
While our $k$-sweep shows stability for $k \in [75, 150]$, a reviewer may reasonably ask whether $k = 75$ was selected post hoc.
We chose $k = 75$ based on prior work conventions for GPT-2 Small circuit analysis and verified stability empirically.

\paragraph{Incomplete ablation comparison.}
The ablation method robustness check (zero vs.\ mean vs.\ resample) was only completed for IOI on GPT-2.
We cannot rule out that the ablation method choice interacts differently with \css{} on other tasks.

\paragraph{IFR--EP correlation.}
On Qwen IOI, the unexpectedly high IFR--EP Jaccard ($0.55$) reduces the effective number of independent methods.
This complicates the interpretation of $K = 3$ consensus on Qwen and suggests that method independence should be verified empirically before applying \css{}.

%======================================================================
\section{Conclusion}
\label{sec:conclusion}

We introduced Circuit Structure Consensus (\css{}), a framework for combining binary circuit discovery outputs across $K \geq 3$ methods.
Rather than claiming that consensus universally improves circuit quality, we characterize \emph{when} it helps: when methods provide independent evidence (low pairwise Jaccard), \css{}$=1.0$ circuits are dramatically more edge-efficient and comparably or more faithful.
When methods agree extensively, consensus degenerates to the most selective method.
We documented a binarization pathology in ACDC that causes spurious degeneration, proposed top-$k$ normalization as a correction, and showed that binary consensus is strictly more stable than score-based ensembling.
These findings provide practitioners with actionable diagnostics for cross-method circuit analysis.

%======================================================================
\bibliographystyle{plain}
\begin{thebibliography}{20}

\bibitem{bhaskar2024finding}
A.~Bhaskar, D.~Friedman, and A.~Ghandeharioun.
\newblock Finding and editing multi-modal neurons in pre-trained transformer models.
\newblock In \emph{Proc.\ NeurIPS}, 2024.

\bibitem{conmy2023towards}
A.~Conmy, A.~N. Mavor-Parker, A.~Lynch, S.~Heimersheim, and A.~Garriga-Alonso.
\newblock Towards automated circuit discovery for mechanistic interpretability.
\newblock In \emph{Proc.\ NeurIPS}, 2023.

\bibitem{elhage2021mathematical}
N.~Elhage, N.~Nanda, C.~Olsson, et al.
\newblock A mathematical framework for transformer circuits.
\newblock \emph{Transformer Circuits Thread}, 2021.

\bibitem{finlayson2021causal}
M.~Finlayson, A.~Mueller, S.~Gehrmann, S.~Shieber, T.~Linzen, and Y.~Belinkov.
\newblock Causal analysis of syntactic agreement mechanisms in neural language models.
\newblock In \emph{Proc.\ ACL}, 2021.

\bibitem{hanna2023how}
M.~Hanna, O.~Liu, and A.~Variengien.
\newblock How does {GPT-2} compute greater-than over the years? {I}nterpreting the failure modes of a trained transformer.
\newblock In \emph{Proc.\ EMNLP}, 2023.

\bibitem{hanna2024have}
M.~Hanna, S.~Balci, and A.~Variengien.
\newblock Have faith in faithfulness: Going beyond circuit overlap in mechanistic interpretability.
\newblock \emph{arXiv:2403.17806}, 2024.

\bibitem{huo2024pie}
X.~Huo, et al.
\newblock {PIE}: Pruning-based iterative ensembling for circuit discovery.
\newblock \emph{arXiv:2604.16889}, 2024.

\bibitem{marks2024sparse}
S.~Marks, M.~Rager, E.~J. Michaud, Y.~Belinkov, D.~Bau, and A.~Mueller.
\newblock Sparse feature circuits: Discovering and editing interpretable causal graphs in language models.
\newblock In \emph{Proc.\ ICML}, 2024.

\bibitem{mondorf2025from}
P.~Mondorf, S.~Ivanova, and B.~Plank.
\newblock From overlap to insight: Comparing methods for circuit discovery.
\newblock \emph{arXiv:2510.06811}, 2025.

\bibitem{parekh2026circus}
B.~Parekh, et al.
\newblock {CIRCUS}: Circuit consensus via threshold diversity.
\newblock 2026.

\bibitem{qwen2025qwen25}
{Qwen Team}.
\newblock Qwen2.5 technical report.
\newblock \emph{arXiv:2412.15115}, 2025.

\bibitem{radford2019language}
A.~Radford, J.~Wu, R.~Child, D.~Luan, D.~Amodei, and I.~Sutskever.
\newblock Language models are unsupervised multitask learners.
\newblock \emph{OpenAI Blog}, 2019.

\bibitem{syed2023attribution}
A.~Syed, C.~Rood, and R.~Fernandes.
\newblock Attribution patching outperforms automated circuit discovery.
\newblock \emph{arXiv:2310.10348}, 2023.

\bibitem{wang2023interpretability}
K.~Wang, A.~Variengien, A.~Conmy, B.~Shlegeris, and J.~Steinhardt.
\newblock Interpretability in the wild: a circuit for indirect object identification in {GPT-2} small.
\newblock In \emph{Proc.\ ICLR}, 2023.

\end{thebibliography}

\end{document}
